% ============================================================
% SCHEDE PG EXTRA — CASA FIESCHI
% @rpg-schema: <https://rpg-schema.org/instances/genova1507/crew-fieschi>
%   fitd:hasOptionalMember
%     <.../pg-caterina-lavagna> ,
%     <.../pg-hamid-rashid> .
% ============================================================

\chapter{Schede PG Extra --- Casa Fieschi}
\index{schede PG!Fieschi!extra}

% ════════════════════════════════════════════════════════════
% PG EXTRA 1 — CATERINA DA LAVAGNA
% @rpg-schema: <https://rpg-schema.org/instances/genova1507/pg-caterina-lavagna>
%   a fitd:Character ;
%   rdfs:label "Caterina da Lavagna" ;
%   fitd:crew <.../crew-fieschi> ;
%   fitd:role "La Spia / L'Infiltrata" ;
%   fitd:action [ fitd:actionName "Survey" ; fitd:dots 4 ] ;
%   fitd:action [ fitd:actionName "Prowl" ; fitd:dots 3 ] ;
%   fitd:action [ fitd:actionName "Sway" ; fitd:dots 3 ] .
% ════════════════════════════════════════════════════════════

\pgheader[FieschiGreen]{Caterina da Lavagna}{La Spia / L'Infiltrata}{Fieschi}
\pgquote{Lavorare per i Fieschi significa essere invisibile. Ho dedicato vent'anni a diventare nessuno.}

\section*{Background \& Motivazione}

\begin{rulebox}{}
  \textbf{Background:} Trentasette anni, nipote di terzo grado di Niccolò,
  abbastanza lontana dalla linea di successione da non contare nulla
  sulla carta. Il conte l'ha reclutata quindici anni fa notando una qualità
  rara: la capacità di essere completamente dimenticabile. Da allora lavora
  come dama di compagnia presso famiglie neutrali — prima mercanti milanesi,
  ora nel palazzo del segretario del governatore francese — osservando,
  ascoltando, e riferendo.

  \medskip\noindent\textcolor{FieschiGreen}{\textbf{Motivazione:}} Caterina
  vuole la rivolta perché vuole smettere. Quindici anni di doppia vita hanno
  un costo che non è visibile nei suoi lineamenti — ma lei lo sente ogni
  mattina. Una Genova libera è una Genova in cui può tornare ad avere
  un nome suo.
\end{rulebox}

\section*{Attributi \& Azioni}

\begin{multicols}{2}

\begin{attributoblock}[FieschiGreen]{INSIGHT --- Mente}
  \azione{Caccia}{Hunt}{2}{Trova informazioni nell'ambiente domestico e amministrativo}
  \azione{Studio}{Study}{2}{Francese, latino, protocollo di corte, stenografia}
  \azione{Osserva}{Survey}{4}{Vent'anni di sopravvivenza basata sull'attenzione}
  \azione{Ingegno}{Tinker}{2}{Codici, messaggi nascosti, comunicazioni sicure}
\end{attributoblock}

\columnbreak

\begin{attributoblock}[FieschiGreen]{PROWESS --- Corpo}
  \azione{Finezza}{Finesse}{3}{Movimenti precisi, presenza fisica calibrata all'invisibilità}
  \azione{Ombra}{Prowl}{3}{Entra ed esce da posti chiusi senza lasciare tracce}
  \azione{Combattimento}{Skirmish}{1}{Sa difendersi — non sa attaccare}
  \azione{Forza bruta}{Wreck}{0}{No}
\end{attributoblock}

\end{multicols}

\begin{attributoblock}[FieschiGreen]{RESOLVE --- Spirito}
  \begin{multicols}{2}
    \azione{Percezione}{Attune}{3}{Sente le tensioni in una stanza prima che diventino visibili}
    \azione{Comando}{Command}{0}{Non comanda — sparisce}
    \azione{Relazioni}{Consort}{2}{Si muove tra servitù, donne di corte, impiegati}
    \azione{Persuasione}{Sway}{3}{Convince con la normalità — è la sua arma principale}
  \end{multicols}
\end{attributoblock}

\medskip
\begin{pgclockbox}{Stress \hfill \clock{9}{0}}
  \small\textit{Traumi: $\square$ Freddo \quad $\square$ Duro \quad $\square$ Ossessionato \quad $\square$ Paranoico \quad $\square$ Irrazionale \quad $\square$ Spezzato}
\end{pgclockbox}

\section*{Abilità Speciale}

% @rpg-schema: <.../pg-caterina-lavagna>
%   fitd:specialAbility [ rdfs:label "La Persona Sbagliata" ] .

\begin{fieschibox}{La Persona Sbagliata}
  Caterina può dichiarare di essere stata scambiata per qualcun altro —
  una dama, una serva, una messaggera — in qualsiasi ambiente sociale.
  Questa identità falsa regge automaticamente per tutta la scena a meno
  che qualcuno non la cerchi specificamente per nome.

  \medskip\noindent\textit{Sempre attiva se non ha già attirato attenzione
  nella scena. Se qualcuno la conosce già, richiede Tiro su Persuasione
  per mantenere la copertura.}
\end{fieschibox}

\section*{Legami}

\begin{table}[h]
\centering\renewcommand{\arraystretch}{1.4}
\begin{tabularx}{\linewidth}{@{}L{2.6cm} X@{}}
  \toprule\rowcolor{FieschiGreen}
  \textcolor{white}{\textbf{PG}} & \textcolor{white}{\textbf{Legame}} \\
  \midrule
  Niccolò & Mi usa da quindici anni. Non mi ha mai ringraziato, ma una volta
           ha detto a qualcuno che ero «la più brava della famiglia».
           Mi ha tenuta in vita quella frase. \\
  \rowcolor{FieschiGreen!8}
  Mons. Gregorio & So del suo accordo con i francesi. Non so ancora se
                  questa informazione mi aiuti o mi metta in pericolo. \\
  Margherita Vedova & Non ci siamo mai parlate. Si capisce che facciamo
                     lo stesso lavoro su piani diversi. \\
  \rowcolor{FieschiGreen!8}
  Ser Elia & L'unico che sa dove sono stanotte. Se non torno, ha istruzioni.
            Non ha chiesto cosa contengano. \\
  \bottomrule
\end{tabularx}
\end{table}

\section*{Equipaggiamento}
\begin{goldlist}
  \item Tre identità complete con documenti diversi — dama, serva, messaggera
  \item Registro mentale delle routine del palazzo del governatore (non scritto)
  \item Piccolo coltello nascosto nell'orlo della gonna
  \item Lettera cifrata pronta da inviare a Niccolò se viene scoperta — contiene tutto ciò che sa
\end{goldlist}

\clearpage

\begin{secretbox}
  \textbf{Segreto:} Stamattina Caterina ha trovato sul suo tavolo una nota
  anonima in francese: \textit{«Sappiamo chi sei. Esci dalla casa stanotte
  o sarai arrestata all'alba.»} Non sa se è vera, se è una trappola per
  farla muovere, o se qualcuno nell'amministrazione francese l'ha coperta
  appositamente per usarla. Non lo ha detto ai Fieschi perché non sa di chi
  fidarsi — inclusi loro.

  \medskip\noindent\textcolor{SecretRed}{\textbf{Agenda:}}
  \textit{Scoprire chi ha scritto quella nota prima della fine della notte.
  Se è vera, fuggire. Se è una trappola, scoprire per chi. Se qualcuno
  la protegge, capire il prezzo.}
\end{secretbox}

\clearpage

% ════════════════════════════════════════════════════════════
% PG EXTRA 2 — HAMID AL-RASHID
% @rpg-schema: <https://rpg-schema.org/instances/genova1507/pg-hamid-rashid>
%   a fitd:Character ;
%   rdfs:label "Hamid al-Rashid" ;
%   fitd:crew <.../crew-fieschi> ;
%   fitd:role "Il Mercante Maghrebino / Il Ponte" ;
%   fitd:action [ fitd:actionName "Hunt" ; fitd:dots 3 ] ;
%   fitd:action [ fitd:actionName "Consort" ; fitd:dots 3 ] ;
%   fitd:action [ fitd:actionName "Sway" ; fitd:dots 3 ] .
% ════════════════════════════════════════════════════════════

\pgheader[FieschiGreen]{Hamid al-Rashid}{Il Mercante Maghrebino / Il Ponte}{Fieschi}
\pgquote{A Tunisi mi chiamano genovese. A Genova mi chiamano straniero. Da nessuna parte mi chiamano per nome.}

\section*{Background \& Motivazione}

\begin{rulebox}{}
  \textbf{Background:} Quarantadue anni, mercante di seta e spezie originario
  di Tunisi, residente a Genova da dodici anni. Ha il \textit{fondaco} nel
  sestiere di Soziglia — tecnicamente sotto protezione Grimaldi, praticamente
  sotto protezione di chiunque gli convenga in quel momento. Il conte Fieschi
  lo ha cercato tre anni fa: Hamid ha accesso alle rotte maghrebine, ai porti
  nordafricani, e — cosa più importante — alla comunità mussulmana che
  commercia con la guarnigione francese senza essere soggetta alle stesse
  restrizioni dei genovesi.

  \medskip\noindent\textcolor{FieschiGreen}{\textbf{Motivazione:}} Hamid
  vuole che la rivolta riesca non per amore di Genova, ma perché i francesi
  hanno tassato il suo \textit{fondaco} al triplo del normale e hanno
  confiscato un carico di seta tunisina «per ragioni di sicurezza» sei mesi
  fa. Quella seta vale seicento ducati. Vuole indietro i seicento ducati.
\end{rulebox}

\section*{Attributi \& Azioni}

\begin{multicols}{2}

\begin{attributoblock}[FieschiGreen]{INSIGHT --- Mente}
  \azione{Caccia}{Hunt}{3}{Reti commerciali, flussi di merci, informatori portuali}
  \azione{Studio}{Study}{2}{Arabo, berbero, italiano, spagnolo, francese commerciale}
  \azione{Osserva}{Survey}{3}{Legge le transazioni e le persone che le fanno}
  \azione{Ingegno}{Tinker}{1}{Logistica, stivaggio, contrabbando tecnico}
\end{attributoblock}

\columnbreak

\begin{attributoblock}[FieschiGreen]{PROWESS --- Corpo}
  \azione{Finezza}{Finesse}{2}{Mani precise da commerciante, sa valutare qualità e difetti}
  \azione{Ombra}{Prowl}{2}{Conosce le vie del porto che non sono sulle mappe ufficiali}
  \azione{Combattimento}{Skirmish}{1}{Sa difendersi — i mercanti a lunga distanza imparano}
  \azione{Forza bruta}{Wreck}{0}{No}
\end{attributoblock}

\end{multicols}

\begin{attributoblock}[FieschiGreen]{RESOLVE --- Spirito}
  \begin{multicols}{2}
    \azione{Percezione}{Attune}{2}{Legge le intenzioni commerciali con precisione}
    \azione{Comando}{Command}{1}{Gestisce i suoi impiegati con autorità pratica}
    \azione{Relazioni}{Consort}{3}{La comunità mercantile straniera a Genova è sua}
    \azione{Persuasione}{Sway}{3}{Negozia in cinque lingue con uguale fluidità}
  \end{multicols}
\end{attributoblock}

\medskip
\begin{pgclockbox}{Stress \hfill \clock{9}{0}}
  \small\textit{Traumi: $\square$ Freddo \quad $\square$ Duro \quad $\square$ Ossessionato \quad $\square$ Paranoico \quad $\square$ Irrazionale \quad $\square$ Spezzato}
\end{pgclockbox}

\section*{Abilità Speciale}

% @rpg-schema: <.../pg-hamid-rashid>
%   fitd:specialAbility [ rdfs:label "Il Fondaco Neutro" ] .

\begin{fieschibox}{Il Fondaco Neutro}
  Il \textit{fondaco} di Hamid è territorio commerciale neutro — non
  genovese, non francese. Qualsiasi incontro organizzato lì avviene
  in uno spazio in cui entrambe le parti abbassano la guardia.
  I tiri sociali condotti nel fondaco partono da Posizione Controllata.

  \medskip\noindent\textit{Sempre attiva quando la scena si svolge nel
  fondaco. Se scoppia violenza nel fondaco, Hamid perde questa abilità
  per il resto della sessione — ha rovinato la neutralità.}
\end{fieschibox}

\section*{Legami}

\begin{table}[h]
\centering\renewcommand{\arraystretch}{1.4}
\begin{tabularx}{\linewidth}{@{}L{2.6cm} X@{}}
  \toprule\rowcolor{FieschiGreen}
  \textcolor{white}{\textbf{PG}} & \textcolor{white}{\textbf{Legame}} \\
  \midrule
  Niccolò & Mi rispetta come si rispetta un fornitore affidabile. Preferisco
           questo all'amicizia — l'amicizia crea aspettative. \\
  \rowcolor{FieschiGreen!8}
  Frà Tomaso & L'unico in questa famiglia che mi ha chiesto come sto.
              Non «come vanno gli affari» — come sto. È una persona strana. \\
  Margherita Vedova & Ha salvato tre dei miei impiegati da un'aggressione
                     l'anno scorso. Non me l'ha mai fatto pesare. \\
  \rowcolor{FieschiGreen!8}
  Ser Elia & Non mi fido dei militari. Ma questo almeno è onesto riguardo
            a cosa è. Lo rispetto per questo. \\
  \bottomrule
\end{tabularx}
\end{table}

\section*{Equipaggiamento}
\begin{goldlist}
  \item Lettere di credito su Tunisi, Ceuta e Barcellona — liquidità ovunque
  \item Rete di cinque informatori tra i commercianti maghrebini che riforniscono la guarnigione francese
  \item Registro delle confische francesi degli ultimi dodici mesi — nomi, date, valori
  \item Passaporto emesso dal sultano hafside di Tunisi — lo rende intoccabile per i diplomatici
\end{goldlist}

\clearpage

\begin{secretbox}
  \textbf{Segreto:} Il carico di seta confiscato dai francesi sei mesi fa
  non è andato perduto — Hamid sa dove si trova: in un magazzino del sestiere
  di Castello, intestato a un funzionario francese che lo rivende
  sistematicamente. Lo stesso funzionario ha accesso ai piani di
  difesa del Castelletto. Hamid ha tenuto questa informazione come
  merce di scambio senza dirlo ai Fieschi — voleva aspettare il momento
  giusto per massimizzarne il valore.

  \medskip\noindent\textcolor{SecretRed}{\textbf{Agenda:}}
  \textit{Usare la notte della rivolta per recuperare la seta e consegnare
  i piani del Castelletto ai Fieschi in cambio di una garanzia commerciale
  permanente. In quest'ordine: prima la seta, poi i piani.}
\end{secretbox}
